\begin{table}[H]
  \centering
  \caption{Matrices de transición de pobreza monetaria por ingreso
  (\% de fila, panel emparejado 1 a 1)}
  \label{tab:matriz_transicion}
  \begin{tabular}{llcc}
    \toprule
    \textbf{Periodo} & \textbf{Estado inicial} & \textbf{No pobre (final)} & \textbf{Pobre (final)} \\
    \midrule
    \multirow{2}{*}{2010 $\rightarrow$ 2013 ($n=8{,}218$)}
      & No pobre & 76.6\% & 23.4\% \\
      & Pobre    & 29.7\% & 70.3\% \\
    \addlinespace
    \multirow{2}{*}{2013 $\rightarrow$ 2016 ($n=6{,}911$)}
      & No pobre & 80.1\% & 19.9\% \\
      & Pobre    & 38.2\% & 61.8\% \\
    \bottomrule
  \end{tabular}
  \begin{minipage}{0.85\textwidth}
    \vspace{4pt}
    \footnotesize \textit{Nota:} $n$ = hogares con emparejamiento 1 a 1;
    excluye 511 (2010--2013) y 1{,}190 (2013--2016) casos de hogares que se
    dividieron entre olas. Fuente: cálculos propios con base en ELCA
    2010, 2013, 2016, generados por \texttt{src/tabla\_matriz\_transicion.py}
    sobre \texttt{outputs/tables/pobreza/transicion\_\{pct,conteo\}\_ola*.csv}
    (\texttt{src/04\_features/build\_pobreza\_desagregaciones.py}).
  \end{minipage}
\end{table}